\section{Design consideration}
\subsection{Responsibilities and cohesion}
\begin{longtable}{@{}p{0.30\textwidth}p{0.64\textwidth}@{}}
\toprule \textbf{Abstraction} & \textbf{Responsibility}\\\midrule\endhead
\texttt{Main}, \texttt{GameCLI} & Construct the program, collect commands, explain invalid choices and display state. The CLI requests actions through the controller.\\
\texttt{GameController} & Coordinate turn phases, staff tasks, arrivals and aggregate results. It owns the game collections and has no console dependency.\\
\texttt{ArrivalSchedule} & Create the scheduled customer type without putting type-dependent patience/pricing rules in the turn loop.\\
\texttt{Customer}, \texttt{VIPCustomer}, \texttt{CriticCustomer} & Own service state and satisfaction; dispatch waiting loss, meal preference and customer discount. Concrete \texttt{Customer} represents a Regular.\\
\texttt{Order} & Own preparation progress, readiness timestamps, locked bill and legal food-state transitions.\\
\texttt{Table} & Hold one occupant until payment or departure; centralise seating.\\
\texttt{Staff} and its four roles & Share the one-task allowance while each subclass exposes its own meaningful operations.\\
\texttt{Menu}, \texttt{MenuItem}, \texttt{ComboMeal} & Describe immutable item prices/preparation and compose the two existing combo components.\\
\texttt{HappyHour}, \texttt{PricingService} & Track activation/duration and calculate the order-time quote with the required discount order.\\
\texttt{PaymentRecord}, \texttt{GameResult} & Preserve completed-payment facts and evaluate the exact victory thresholds.\\
\texttt{TurnPhase}, \texttt{CustomerStatus}, \texttt{OrderStatus} & Hold named String constants. These are not Java enums.\\
\bottomrule
\end{longtable}

\subsection{Principles reflected in this code}
\textbf{SRP.} The CLI handles interaction, the controller schedules the turn,
Order guards preparation/serving/payment transitions, and PricingService quotes
the bill. These responsibilities can be discussed and tested separately.
The small forecast shown by GameCLI is presentation logic with timing
assumptions; keeping it beside the display is a simplicity trade-off.

\textbf{OCP.} Customer references dispatch \texttt{getWaitingLoss()},
\texttt{chooseMeal()} and \texttt{discountedBill()} to the applicable subclass.
Adding VIP/Critic behaviour did require setup and a pricing hook in Stage~2;
the claim is reduced coupling, not zero changes to existing code.
Adding Host and the pricing policy did not require rewriting Chef or Cashier.

\textbf{LSP.} Regular, VIP and Critic customers follow the same seat/order/serve/
eat/pay flow. Their overrides change documented preferences, patience and
pricing while preserving that service contract. ComboMeal supports the same
MenuItem getters used by ordering and cooking.

\textbf{ISP.} There is no interface hierarchy in this small implementation.
Staff shares only task availability; it does not force Cashier to implement
cooking or Chef to implement seating. This is a focused API choice, not a
claim that an unused interface pattern has been implemented.

\textbf{DIP.} Waiter receives a PricingService through its constructor and
GameCLI receives its controller and Scanner. These explicit dependencies aid
construction and testing. Both still depend on concrete classes; full
interface-based dependency inversion is deliberately not claimed.

\subsection{Encapsulation, composition and alternatives}
Private fields and domain operations guard state. A failed action leaves the
controller in the same phase, and Staff completes a task only after a valid
action. Order and Customer timestamps prevent same-turn serving and early payment, respectively.
PaymentRecord is immutable, so later events cannot change recorded revenue or
paid satisfaction. Returned ArrayLists are container copies containing the same
domain objects; this is a deliberately limited boundary for a six-customer CLI.

Customer has at most one Order. GameController owns its staff, tables and
payment records. ComboMeal holds references to two existing immutable MenuItem
components; this is aggregation of shared menu objects rather than exclusive
ownership of them. The waiting queue is derived from customer state in arrival
order, avoiding a second membership collection that could become inconsistent.

An alternative was one large controller with customer-type conditionals,
parallel order arrays and discount branches. It would reduce the number of
files, but entangle service transitions, pricing and display. The chosen model
keeps those rules in the objects that own the affected state. A more generic
engine with interfaces for every phase was also unnecessary for 18 turns and
six customers. Named String constants are familiar to the course, though enums
would give stronger compile-time state checking in a larger project.
